\documentclass[10pt,a4paper]{amsart}
\usepackage[margin=19mm]{geometry}
\usepackage{amsmath,amssymb,mathtools}
\usepackage[scr=rsfs,cal=euler]{mathalfa}
\usepackage{xfrac}
\usepackage[unicode,hidelinks,bookmarks=false]{hyperref}
\newcommand{\ie}{i.e.\ }
\newcommand{\cf}{cf.\ }
\newcommand{\resp}{respectively\ }
\newcommand{\Subsection}{\subsection}
\providecommand{\nobreakdash}{\nobreak\hbox{-}}
\setlength{\emergencystretch}{2em}
\multlinegap0pt
\usepackage{bm}
\usepackage{bbm}
\usepackage{dsfont}
\usepackage{xspace}
\usepackage{cancel}
\usepackage{mathtools}
\usepackage{amsthm}
\makeatletter
\@ifundefined{theorem}{\newtheorem{theorem}{Theorem}}{}
\@ifundefined{hypothesis}{\newtheorem{hypothesis}{Hypothesis}}{}
\@ifundefined{lemma}{\newtheorem{lemma}[theorem]{Lemma}}{}
\makeatother
\newcommand{\C}{\mathcal C}
\title[Coloring graphs as complete graph invariants]{Coloring graphs as complete graph invariants}
\author{Shamil Asgarli, Sara Krehbiel, Howard W. Levinson}
\date{Source: arXiv:2504.20978v3 (CC BY 4.0)}
\begin{document}
\maketitle

\noindent\emph{Source and licence.} Coloring graphs as complete graph invariants. Version arXiv:2504.20978v3 (CC BY 4.0), distributed under the Creative Commons Attribution 4.0 International licence, as recorded in the arXiv licence metadata for this exact version and shown on the abstract page https://arxiv.org/abs/2504.20978v3. Full rights notice: licenses/arxiv-2504.20978-CC-BY-4.0.txt.\par\smallskip

\noindent\textit{Editorial scope.} Counted unit: the Lemma whose source label is \texttt{lem:squares-in-linked-graph} in the article (source0.tex lines 1652--1663, source label \texttt{lem:squares-in-linked-graph}), delivered together with the article's own complete original proof of it (source0.tex lines 1664--1684, one \texttt{proof} environment of 21 lines). No mathematics is written, repaired or re-derived: statement and proof are the article's own text line for line. Required context retained uncounted: lem:subgraphs (lines 879--881); hyp:running (lines 1091--1092). Every definition, display and label of those lines is retained unchanged. Omitted from this package and not redistributed: 1--878, 882--1090, 1093--1651, 1685--2065. Nothing inside a retained excerpt is omitted or altered: every retained body is delivered exactly as the article's lines are written, so no omission receipt of any kind accompanies this package. Citations: the retained excerpts carry no citation command at all, so no bibliography entry of the article is transcribed and no reference is redistributed. Reference adaptation: none was needed and none was made; every reference inside the delivered unit resolves inside it, so no unresolved reference or citation is produced. Printed slips kept literal: no printed word, symbol, hypothesis or number of the counted statement or of its proof was altered. Layout and class: the article's own class is \texttt{amsart} with packages including fullpage, colonequals, amsrefs, enumitem, amssymb, stmaryrd, mathtools, xcolor, tikz, hyperref, tikz-cd, adjustbox, ulem; the wrapper is the lane's pinned \texttt{amsart} editorial wrapper at 10pt on A4, which restates the theorem-like environments the retained text needs and adds the article's own macro definitions that the retained text needs (\texttt{\textbackslash C}), with extra wrapper packages (bm, bbm, dsfont, xspace, cancel, mathtools). The delivered numbering is the unit's own. Assets: no retained span contains a figure environment, a graphics-inclusion command, a PDF-page inclusion, a table or a TikZ picture; no image file is copied into this package, the package declares no asset dependency and no image file is redistributed with it. Licence: version arXiv:2504.20978v3 of this article is distributed under the Creative Commons Attribution 4.0 International licence, as recorded in the arXiv licence metadata for this exact version and shown on the abstract page https://arxiv.org/abs/2504.20978v3. A full rights notice is delivered at licenses/arxiv-2504.20978-CC-BY-4.0.txt. Characters: every retained excerpt is pure ASCII, so the delivered engine, XeLaTeX with its default Latin Modern text font, reports no missing character.
\begin{lemma}\label{lem:subgraphs}
    If $\C\cong \C_k(G)\cong \C_{\chi'}(G')$ where $k>\chi(G)$ and $\chi'=\chi(G')$, then $G'$ contains $G$ as a subgraph. Moreover, if $G$ consists of connected components $H_1,\dots,H_d$, then $G'$ contains the connected components of $G$ as vertex-disjoint induced subgraphs. That is, there exists an injection $\phi\colon V(G)\to V(G')$ such that for any $i\in[d]$ and $u,v\in H_i$, we have $uv\in E(G)$ if and only if $\phi(u)\phi(v)\in E(G')$. 
\end{lemma}

\begin{hypothesis}\label{hyp:running} Suppose $\C\cong \C_k(G)\cong \C_{\chi'}(G')$ where $k>\chi(G)$ and $\chi'=\chi(G')$. This pair of isomorphisms, $\C\cong \C_k(G)$ and $\C_k(G)\cong \C_{\chi'}(G')$, will be fixed, as will the underlying colorings of $G$ and $G'$ corresponding to each vertex in $\C$. Assume $H_1, \dots, H_d$ are the connected components of $G$. If $\alpha\in V(\C)$ is an abstract link vertex, then $\alpha_G$ and $\alpha_{G'}$ will be the corresponding link $k$-coloring of $G$ and $\chi'$-coloring of $G'$, respectively.  By Lemma~\ref{lem:subgraphs}, $G'$ contains subgraphs $H'_1, \dots, H'_d$ such that $H'_i\cong H_i$ for each $i\in [d]$. We also recall the embedding $\phi\colon V(G)\to V(G')$ from Lemma~\ref{lem:subgraphs}. 
\end{hypothesis}

\begin{lemma}\label{lem:squares-in-linked-graph} Assume Hypothesis~\ref{hyp:running}.  Let $\alpha, \alpha_1, \alpha_2 \in V(\mathcal{L}'_{[\alpha]})$ be distinct vertices such that $\alpha$ is adjacent to both $\alpha_1$ and $\alpha_2$. 
\begin{enumerate}
    \item If $\alpha\alpha_1$ has label $c_1Q_{1}^ic_3$ and $\alpha\alpha_2$ has label $c_2 Q_{2}^ic_4$, where $c_1,c_2,c_3$, and $c_4$ are all distinct, and $Q_{1}^i$ and $Q_2^{i}$ are disjoint sets, then there is a unique induced $4$-cycle in $\mathcal{L}'_{[\alpha]}$ that contains these two edges and its labels are given by:
     \[\alpha \xleftrightarrow{c_1Q^i_{1}c_3}  \alpha_1\xleftrightarrow{c_2Q^i_{2}c_4} \beta\xleftrightarrow{c_3Q^i_{1}c_1}\alpha_2\xleftrightarrow{c_4Q^i_{2}c_2}  \alpha.
\]
    \item If $\alpha\alpha_1$ has label $c_1Q_{1}^ic_3$ and $\alpha\alpha_2$ has label $c_1 Q_{1}^ic_4$, where $c_1,c_3$, and $c_4$ are all distinct, then there is a unique $3$-cycle in $\mathcal{L}'_{[\alpha]}$ that contains these two edges and its labels are given by: 
    \[\alpha \xleftrightarrow{c_1Q^i_{1}c_3}  \alpha_1\xleftrightarrow{c_3Q^i_{1}c_4} \alpha_2\xleftrightarrow{c_4Q^i_{1}c_1}  \alpha.
\]
    \item If $\alpha\alpha_1$ has label $c_1Q_{1}^ic_3$ and $\alpha\alpha_2$ has label $c_2 Q_{2}^jc_4$ for $i\neq j$, where necessarily $c_1\neq c_3$ and $c_2\neq c_4$, and $Q_1^i$ and $Q_2^j$ are disjoint sets, then there is a unique induced $4$-cycle in $\mathcal{L}'_{[\alpha]}$ that contains these two edges and its labels are given by: \[\alpha \xleftrightarrow{c_1Q^i_{1}c_3}  \alpha_1\xleftrightarrow{c_2Q^j_{2}c_4} \beta\xleftrightarrow{c_3Q^i_{1}c_1}\alpha_2\xleftrightarrow{c_4Q^j_{2}c_2}  \alpha.
\]
\end{enumerate}
\end{lemma}

\begin{proof}
We prove each statement independently below.
\begin{enumerate}
    \item We are given $\alpha_{G'}(Q^i_1)=c_1, \alpha_{G'}(Q^i_2)=c_2$, with edges $\alpha \xleftrightarrow{c_1Q^i_1c_3} \alpha_1$ and $\alpha \xleftrightarrow{c_2Q^i_2c_4} \alpha_2$. Thus, $(\alpha_1)_{G'}(Q^i_1)=c_3, (\alpha_1)_{G'}(Q^i_2)=c_2$ and $(\alpha_2)_{G'}(Q^i_1)=c_1, (\alpha_2)_{G'}(Q^i_2)=c_4$.
    Since $Q^i_1, Q^i_2$ are disjoint parts and $c_1,c_2,c_3,c_4$ are distinct, the recoloring operations $Q^i_1\colon c_1 \to c_3$ and $Q^i_2\colon c_2 \to c_4$ are independent. Let $\beta$ be the coloring where $\beta_{G'}(Q^i_1)=c_3$ and $\beta_{G'}(Q^i_2)=c_4$ (other parts as in $\alpha$). This $\beta$ is proper as $c_3 \neq c_4$. The edges $\alpha_1 \xleftrightarrow{c_2Q^i_2c_4} \beta$ and $\beta \xleftrightarrow{c_3Q^i_1c_1} \alpha_2$ complete the stated unique induced $4$-cycle $\alpha \xleftrightarrow{c_1Q^i_1c_3} \alpha_1 \xleftrightarrow{c_2Q^i_2c_4} \beta \xleftrightarrow{c_3Q^i_1c_1} \alpha_2 \xleftrightarrow{c_4Q^i_2c_2} \alpha$.
    \item  We are given $\alpha_{G'}(Q^i_1)=c_1$, with edges $\alpha \xleftrightarrow{c_1Q^i_1c_3} \alpha_1$ and $\alpha \xleftrightarrow{c_1Q^i_1c_4} \alpha_2$. Thus, $(\alpha_1)_{G'}(Q^i_1)=c_3$ and $(\alpha_2)_{G'}(Q^i_1)=c_4$, while all other parts are colored as in $\alpha$. Since $c_3 \neq c_4$, $\alpha_1$ and $\alpha_2$ differ only on $Q^i_1$ and are thus adjacent via the edge $\alpha_1 \xleftrightarrow{c_3Q^i_1c_4} \alpha_2$. Together with the edge $\alpha_2 \xleftrightarrow{c_4Q^i_1c_1} \alpha$ (returning $Q^i_1$ to its color in $\alpha$), this forms the unique induced $3$-cycle $\alpha \xleftrightarrow{c_1Q^i_1c_3} \alpha_1 \xleftrightarrow{c_3Q^i_1c_4} \alpha_2 \xleftrightarrow{c_4Q^i_1c_1} \alpha$.
    \item We are given $\alpha_{G'}(Q^i_1)=c_1, \alpha_{G'}(Q^j_2)=c_2$, with initial edges $\alpha \xleftrightarrow{c_1Q^i_1c_3} \alpha_1$ and $\alpha \xleftrightarrow{c_2Q^j_2c_4} \alpha_2$. Here $Q^i_1$ and $Q^j_2$ are disjoint vertex sets, $c_1 \neq c_3$, and $c_2 \neq c_4$. The recolorings $Q^i_1\colon c_1 \to c_3$ and $Q^j_2\colon c_2 \to c_4$ are independent. If $c_3 \neq c_4$, the reasoning from part (1) applies directly: the independent operations with distinct target colors define the unique induced 4-cycle and its labeling as stated.
    
   Consider the case $c_3 = c_4 = c$. We have $\alpha_{G'} \equiv (Q^i_1\colon c_1, Q^j_2\colon c_2)$, $(\alpha_1)_{G'} \equiv (Q^i_1\colon c, Q^j_2\colon c_2)$, and $(\alpha_2)_{G'} \equiv (Q^i_1\colon c_1, Q^j_2\colon c)$. These vertices complete to a 4-cycle in $\mathcal{L}'_{[\alpha]}$ because the matching $4$-cycle exists in $\mathcal{L}_{[\alpha]}$. Let $\beta$ be the fourth vertex in this 4-cycle $(\alpha, \alpha_1, \beta, \alpha_2,\alpha)$.
    The $\alpha$-avoiding path from $\alpha_1$ to $\alpha_2$ is $\alpha_1 \to \beta \to \alpha_2$. To change from $(\alpha_1)_{G'} \equiv (Q^i_1\colon c, Q^j_2\colon c_2)$ to $(\alpha_2)_{G'} \equiv (Q^i_1\colon c_1, Q^j_2\colon c)$, part $Q^i_1$ must change from $c$ to $c_1$, and part $Q^j_2$ must change from $c_2$ to $c$.
    These two changes must occur over the two edges $(\alpha_1, \beta)$ and $(\beta, \alpha_2)$. There are two ways to assign these changes to the edges:
    \begin{enumerate}
        \item Path $T_1$: $\alpha_1 \xleftrightarrow{cQ^i_1c_1} \beta \xleftrightarrow{c_2Q^j_2c} \alpha_2$.
        Here, $\beta$ would have $Q^i_1$ colored $c_1$ and $Q^j_2$ colored $c_2$. Thus, $\beta_{G'} \equiv (Q^i_1\colon c_1, Q^j_2\colon c_2)$, which is $\alpha$. This is a contradiction, because $\beta$ is distinct from $\alpha$.
        \item Path $T_2$: $\alpha_1 \xleftrightarrow{c_2Q^j_2c} \beta \xleftrightarrow{cQ^i_1c_1} \alpha_2$.
        Here, from $(\alpha_1)_{G'} \equiv (Q^i_1\colon c, Q^j_2\colon c_2)$, recoloring $Q^j_2\colon c_2 \to c$ yields $\beta_{G'} \equiv (Q^i_1\colon c, Q^j_2\colon c)$. This path $T_2$ avoids $\alpha$, so this is the correct labeling.
    \end{enumerate}
    Therefore, $T_2$ provides the unique labeling for the $\alpha$-avoiding path segment $\alpha_1 \to \beta \to \alpha_2$. Reinstating $c$ as $c_3$ (for $Q^i_1$'s target) and $c_4$ (for $Q^j_2$'s target, which is also $c_3$ in this subcase), the labels become $\alpha_1 \xleftrightarrow{c_2Q^j_2c_4} \beta$ and $\beta \xleftrightarrow{c_3Q^i_1c_1} \alpha_2$.
    This completes the 4-cycle as stated in the lemma: $\alpha \xleftrightarrow{c_1Q^i_1c_3} \alpha_1 \xleftrightarrow{c_2Q^j_2c_4} \beta \xleftrightarrow{c_3Q^i_1c_1} \alpha_2 \xleftrightarrow{c_4Q^j_2c_2} \alpha$. In retrospect, the properness of $\beta$, ensured by the existence of this 4-cycle structure, implies that if $c_3=c_4$, no edge exists between $Q^i_1$ and $Q^j_2$. \qedhere
\end{enumerate}
\end{proof}

\end{document}
